\section{Introduction}
\label{sec:intro}

Data pipelines repeatedly collect, transform, and deliver records for analysis and operational decisions. A pipeline may load flight records each day, invoice lines each hour, or inspection outcomes each month. Declarative quality rules guard these batches by checking properties such as completeness, uniqueness, numerical plausibility, and membership in a valid category domain. Systems including Deequ~\cite{Schelter2018,Schelter2019}, TensorFlow Data Validation (TFDV)~\cite{Breck2019,Caveness2020}, and Data Sentinel~\cite{Swami2020} make such checks repeatable and support profiling or constraint management. Their effectiveness nevertheless depends on how validation expectations are established and maintained as the data change.

Three connected difficulties motivate this study. First, a fixed expectation can become outdated even when the pipeline is working correctly. Seasonal demand, congestion, and revised reporting procedures can alter valid values, causing rules learned from an earlier sample to reject later records. In the flight-data experiment, a mean-plus-or-minus-three-standard-deviations rule calibrated on five warm-up batches flags approximately 4.6\% of valid records at 2\% injected drift and 7.6\% at 30\% drift, despite an unchanged 1\% background error rate. A growing alarm rate can therefore reflect legitimate evolution rather than deteriorating quality.

Second, learning from recent batches can create the opposite problem. An adaptive range may follow valid changes in location and scale, but it may also widen around recurring numerical corruption. Likewise, repeated observation of a previously unseen country name does not establish that the name is valid: the value may be a persistent misspelling. A monitor must balance responsiveness to change against preservation of the constraints that define errors. Reducing the number of alarms alone is not evidence that this balance has been achieved.

Third, validation consumes resources. Executing every check on every batch can repeat scans and spend effort on columns that show little evidence of change. Risk-based validation proposes directing effort toward the probability and consequence of data problems~\cite{FoidlFelderer2019}. Its practical value, however, depends on whether selection preserves detection and whether skipped checks represent a substantial share of total monitoring work. Profiling and drift estimation may remain necessary even when some checks are omitted.

These difficulties also affect the interpretation of an alarm. An operator needs to know whether a batch differs from history because valid data have changed, because errors have increased, or because both occurred. We refer to legitimate distribution change as \emph{drift} and increased erroneous-record prevalence as \emph{degradation}. These properties are distinct but can coexist. A distribution-change signal alone cannot explain which condition holds; a rule-failure signal alone can mistake outdated expectations for poor quality. The framework therefore combines the two signals while retaining separate measurements of record-level detection and batch-level classification.

Relevant research supplies several of the necessary building blocks. History-aware validation learns expectations from recurring pipeline executions~\cite{Redyuk2021,Tu2023,Shankar2023}. Drift research provides methods for comparing distributions~\cite{Gama2014,Lu2019,Rabanser2019}. LLMs can translate schema descriptions, examples, and requirements into candidate quality conditions~\cite{Glock2026,Xie2025,Abughazala2026,Chen2026}. These developments motivate an integrated evaluation rather than an assumption that combining them necessarily improves performance. In particular, useful evidence must separate the contribution of threshold adaptation from that of drift gating, distinguish domain knowledge from historical frequency, and compare selective execution with simpler schedules at a similar check budget.

This article develops a configurable batch monitor to investigate those questions within one framework. It profiles each batch, constructs static and statistical checks, admits LLM and dictionary-only model proposals after warm-up validation, adapts numerical ranges, and selects checks through a per-column risk score. A gated distribution comparison and a failure-rate signal then assign one of four batch labels. Category-guard policies are evaluated separately to determine which forms of domain evidence can constrain the admission of recurring new values.

The evaluation addresses three research questions. RQ1 asks whether adaptive thresholds reduce false positives under legitimate drift without an unacceptable recall loss, and whether the combined signals distinguish the four batch conditions. RQ2 asks whether admitted LLM rules add coverage of domain-impossible values and whether domain checks improve category admission. RQ3 asks whether risk-based check selection reduces execution cost while maintaining detection after degradation begins. These objectives are evaluated individually, so that success of one component is not treated as success of every configuration.

Three public datasets cover transport, retail, and public health. Seventeen controlled scenarios provide audited injected-error labels, while disjoint development and test record pools separate method development from the paired test evaluation. The internally frozen protocol specifies the margins, settings, and ten test seeds. Comparators include pooled fixed, rolling, robust, and exponentially weighted thresholds; an AVH-inspired historical-statistics baseline~\cite{Tu2023}; and configured Deequ and TFDV runs. Two dirty/clean benchmark pairs~\cite{Mahdavi2019} and a chronological Chicago replay examine transfer beyond shuffled controlled batches.

The article makes the following contributions.
\begin{enumerate}
\item An integrated monitoring loop that combines gated drift evidence and rule-failure evidence to assign four operational batch labels, with warm-up drift calibration based on leave-one-out scoring and resampling (Section~\ref{sec:method}).
\item A paired evaluation protocol with disjoint development and test records, prespecified non-inferiority margins, seed-cluster uncertainty estimates, and manifests recording code and configuration fingerprints (Section~\ref{sec:design}).
\item Dataset- and scenario-specific evidence on numerical and categorical adaptation, including reductions in drift-related false positives and losses in detection when persistent errors enter learned expectations (Section~\ref{sec:results}).
\item Evaluation of LLM-proposed rules and their admission procedure, including complementary statistical coverage, deliberately faulty candidates, dictionary-only model comparisons, and category-domain checks.
\item A cost analysis that distinguishes executed checks, estimated bytes touched, and measured monitoring phases, with matched scheduling baselines and a profile-reuse configuration.
\end{enumerate}

The results support conditional benefits. Four-class macro-F1 reaches 0.966, 0.955, and 0.947 on the three datasets, but the same drift gate with fixed thresholds gives identical Chicago results and no detected disadvantage on Retail. Threshold adaptation meets the primary false-positive and recall criteria in six of nine flight scenarios and nine of nine Chicago scenarios; Chicago's small gain comes from category policy. It does not meet the false-positive criterion on Retail. Fixed statistical rules combined with admitted LLM rules reach semantic recall of 0.996--1.00, whereas recurring invalid codes remain a weakness of adaptive combinations. Risk-based execution meets the point-estimate usefulness criterion in eleven of eighteen degradation/incident cases, with small absolute savings and little timing benefit after profile reuse. These findings identify both the contribution and the operating limits of the proposed framework.

The article is organised as follows. Section~\ref{sec:related} reviews related work. Section~\ref{sec:problem} defines the setting and hypotheses. Section~\ref{sec:method} explains the monitor, and Section~\ref{sec:design} describes the evaluation. Section~\ref{sec:results} reports the hypothesis results, ablation, external checks, and sensitivity analysis. Sections~\ref{sec:discussion} and~\ref{sec:conclusion} discuss their implications and conclude.
